\documentclass[11pt,letterpaper]{article}
\usepackage[margin=0.7in,headheight=15pt]{geometry}
\usepackage[T1]{fontenc}\usepackage{lmodern,amsmath,amssymb,fancyhdr,parskip,hyperref}
\hypersetup{colorlinks=true,urlcolor=blue}
\pdfinfoomitdate=1\pdftrailerid{}\pdfsuppressptexinfo=15
\pagestyle{fancy}\fancyhf{}\lhead{Week 79 / Numbered-ball detectives}\rhead{Adult guide}\cfoot{Bellingham Math Circle / Week 79 / facilitator / \thepage}
\setlength{\parskip}{7pt}\setlength{\parindent}{0pt}
\begin{document}
\section*{Mathematical overview}
Under the smallest-card rule, after completed day $n$, $A_n=\{n+1,\ldots,2n\}$ and $B_n=\{1,\ldots,n\}$. Both counts are $n$. Every fixed positive label $k$ is introduced on day $\lceil k/2\rceil$ and is in $B$ permanently by day $k$. Under the \emph{eventual-membership} interpretation, the limiting sets are empty $A$ and all positive labels in $B$. There is no last finite day.

If instead we transfer the larger of the two new labels, after day $n$ the odd labels through $2n-1$ remain in $A$ and the even labels through $2n$ are in $B$. The finite counts still tie. Count histories do not determine limiting membership.

If the mover may choose any card in $A$ each day, a set of $n$ labels can be $B$ after day $n$ exactly when, for every day $d\le n$, at least $d$ of its labels are at most $2d$: label $k$ cannot move before it enters on day $\lceil k/2\rceil$, and when the condition holds, moving the smallest wanted label in $A$ each day always works.

Experiments establish small cases. Induction and a named label's finishing day explain all cases. Children are meant to distinguish ``each label eventually'' from ``all labels by one day.'' Band map: Grades 2--5 (pages 1--3) enact the rule, follow named cards, compare the two rules' counts and cards, and choose moves to reach a target B (Problems 1--4). Grades 4--5 (page 4) decide where every card settles under each rule and whether a finite day shows it (Problem 5). K--1 has no pages; page 2 suggests an adult-run route. This is a new teaching adaptation, not a claim that actual infinitely many days have occurred.

The title recalls the organizer's numbered-ball process. The pages use numbered cards; this guide calls the number on a card its label.
\newpage
\section*{Preparation and a first route}
For each pair: cards 1--24 and either two shallow bowls OR the supplied A/B mat as the single shared state; place unused cards in consecutive pairs outside both bowls. Add one extra card sheet and mat per table, so Problem 3 can run both rules side by side. One adult anchors each participating table. For the three-child upper table, rotate feeder, mover and label detective. Adults record a few common observations; do not require children to track counts and every label at once.

Reading: short rules read aloud if needed. Arithmetic: compare labels, count to 24, add two and remove one. Reasoning: track one object across several changing arrangements. The limit discussion is optional; do not mistake arithmetic speed for readiness.

After the usual five-minute movement break, allow two minutes to handle/sort cards. Launch together for three minutes: demonstrate adding a pair, locating the \emph{smallest label anywhere in A}, and moving it to B. The printed before/intermediate/after visual supports the convention. Show a legal finite move, not the limiting answer.

First route: 15 minutes enacting the first six days and predicting arrangements (Problem 1); twelve days is an optional continuation; 10 minutes with children adopting labels (Problem 2); 10 minutes on the changed rule, both rules side by side (Problem 3); five minutes of shared explanations. Leave at least five minutes for questions or replay. A table that finishes early takes Problem 4 (chosen moves, about 10 minutes); otherwise it opens a return visit. Grades 4--5 may go on to Problem 5 instead. Use one mat per rule repeatedly; bring in later questions when the table is ready. Draft timing and physical handling are untested.

K--1 route (no K--1 pages; a prediction from the review, untested): an adult reads the rule and adds each day's two cards, and the child finds and moves A's smallest card. Run Problem 1's six days with cards 1--12 and count both bowls at the end of each day. Then run card 1 under both rules, side by side: after day 1 it is in B under the smallest-card rule and in A under the changed rule, and it stays there. Neither needs reading; the adult keeps any record. Problem 2's prediction for card 20 and Problem 5 are beyond this route.
\newpage
\section*{Facilitator reasoning and hints}
Induction: assuming $A_n=\{n+1,\ldots,2n\}$, adding $2n+1,2n+2$ and transferring $n+1$ gives $A_{n+1}=\{n+2,\ldots,2n+2\}$. The transferred label appends to $B$. The empty day 0 is the base case. Problem 1: adding two and moving one increases each bowl's count by one, so both have $n$ cards after day $n$. Day 6 gives $A=\{7,\ldots,12\}$ and $B=\{1,\ldots,6\}$. Day 10 gives $A=\{11,\ldots,20\}$ and $B=\{1,\ldots,10\}$.

Problem 2: label 3 enters on day 2 and moves on day 3; 6 enters on day 3 and moves on day 6; 10 enters on day 5 and moves on day 10; 20 enters on day 10 and moves on day 20. For any $k$, entry is $\lceil k/2\rceil$ and transfer is day $k$, so a card's wait in A grows with its label: card 20 waits ten days. That growing wait is what Problem 5 turns on. When a child adopts label $k$, ask only when that card enters and when it leaves A. For example label 8 enters on day 4, remains in A through day 7 and moves on day 8. An unseen label can be discussed with a placeholder; the limited card supply does not stop the mathematical rule.

Problem 3: run the two rules side by side rather than from memory. At a table with two pairs, one pair reruns the smallest-card rule while the other runs the changed rule, for the same number of days; at a three-child Grades 4--5 table, the third child runs the second rule on the extra card sheet and mat. Then compare the two actual sets. After day 1, card 1 is in B under the smallest-card rule and in A under the changed rule; the counts are equal every day.

Problem 4 (chosen moves): the first two rows can be reached; the third cannot. Row 1: move the smaller new card each day (1, 3, 5, 7, 9, 11). Row 2: for example 2, 3, 5, 7, 8, 11 on days 1--6. Row 3: by the end of day 4, four cards have moved and only cards 1--8 have entered, but only three row-3 cards (1, 4, 6) are among 1--8, so some card outside the row is in B by day 4, and cards in B stay. A run usually shows this as a stuck day: after moving 1, 4 and 6, day 4's A is 2, 3, 5, 7, 8. Ask ``Could a different order work?'' only after a child meets the stuck day. This is the first problem where children, not a rule, make the moves, and the first ``no'' that must cover every way of playing; it uses Problem 2's entry day. For adults: the counting condition in the overview is exact, and the package checker confirms it by trying every sequence of moves through day 7. The reachable sets number 2, 5, 14, 42, 132, 429, 1430 for days 1--7, the Catalan numbers.

If the table says A and B always hold the same number, agree that the finite count claim is correct. Ask what happens to one \emph{fixed} label. Problem 5, smallest-card rule: inspect a named finite list of labels, leaving every A/B card where it is. Inspect labels 1--6, then 1--12, regardless of location. Any fixed finite list eventually has all its labels in B. At every positive finite day $N$, label $N+1$ is still in A. Day 0 is empty because no labels have entered, and is not a day when all positive labels are in B. Hence no finite day finishes every label.

Problem 5, changed rule: the larger new card is $2n$, so it moves immediately to B. The smaller $2n-1$ never moves. So the odd labels settle in A and the even labels in B. On every finite day each card present is already in its settled bowl, so ``it already looks like that'' is right about the cards present. The whole arrangement never occurs, because labels above $2N$ have not arrived by day $N$. Contrast the smallest-card rule, where cards $N+1$ to $2N$ sit in A, the bowl they will leave. Do not print the odd/even answer on their working board.
\newpage
\section*{Limits, encore, sources and observation record}
This is a two-card version of the Ross--Littlewood paradox (the balls-and-vase problem), named for J.~E. Littlewood and Sheldon Ross. Adult analysis: the indicators $a_n(k)=1_{A_n}(k)$ tend pointwise to 0, while $\sum_k a_n(k)=n\to\infty$. Thus $\sum_k\lim_n a_n(k)=0$ differs from $\lim_n\sum_k a_n(k)$. Convergence is not uniform: $\sup_k a_n(k)=1$ for every $n\ge1$ (day 0 is empty). Infinite counting measure and the lack of a summable majorant matter; nonuniformity alone is not a universal explanation for every failed interchange. For B, membership is monotone and its count tends to infinity. None of this establishes an extra ``midnight state'' without selecting the membership interpretation.

Encore: children invent a transfer rule while keeping the two-per-day/one-transfer convention. Investigate labels that eventually move, labels that never move, and labels whose locations never settle if backward transfers are newly permitted. State any changed rules explicitly. Connect later to the midpoint process in Week 80, or existing Week 50 staircase limits. Those are different questions about limits.

Source: the organizer's numbered-ball process and changed-rule proposal from the October 9 brief; finite formulas, comparison and proof are independently derived here. Problem 4's chosen-move targets come from the October 17 review card (fix 3), checked by exhaustive search. Pedagogy: Rozhkovskaya, \emph{Math Circles for Elementary School Students}, Lesson 8, ``At the lesson,'' item 2 (EPUB part0018): instructors reduced slow copying. Our adult-scribed record is an adaptation.

Unpiloted. Print US Letter, single-sided, Actual Size. Physical fit, card handling and classroom procedure have not been rehearsed. Record afterward: tasks actually used; child actions and words; confusion requiring rescue; questions they wanted to continue. Separate observations from explanations or hypotheses.
\end{document}
